\documentclass[11pt,a4paper]{article}
\usepackage[margin=1in]{geometry}
\usepackage{amsmath,amssymb,amsthm}
\usepackage{algorithm}
\usepackage{algpseudocode}
\usepackage{booktabs}
\usepackage{hyperref}

\theoremstyle{definition}
\newtheorem{definition}{\textbf{Definition}}[section]
\newtheorem{theorem}{\textbf{Theorem}}[section]
\newtheorem{lemma}[theorem]{\textbf{Lemma}}
\newtheorem{remark}{\textbf{Remark}}[section]

\title{\textbf{Weight-Decomposed Low-Rank Adaptation (DoRA) and Rank-Stabilized LoRA (rsLoRA)}}
\author{\textbf{Team Scaibu} \\ \small Email: \texttt{jaydeep.9664920749@gmail.com} \\ \small \textbf{Architecture and Core Engine Team}}
\date{\textbf{October 2026}}

\begin{document}
\maketitle

\section{Definitions and Cognitive Mental Models}

\subsection{Formal Mathematical Definition}
\textbf{Definition (Weight Decomposition and Directional Low-Rank Tuning):}
Standard Low-Rank Adaptation (LoRA) parametrizes weight changes via an additive matrix delta $\Delta W = s \cdot BA$, modifying magnitude and directional components simultaneously in a coupled fashion. Empirical analysis reveals that full fine-tuning (FT) exhibits distinct variations: magnitude changes tend to have low correlation with directional changes, whereas standard LoRA shows an artificially high positive correlation between magnitude and directional updates.

\textbf{Weight-Decomposed Low-Rank Adaptation (DoRA)} disassembles any weight matrix $W \in \mathbb{R}^{d \times k}$ into its magnitude vector $m \in \mathbb{R}^d$ and normalized directional matrix $V \in \mathbb{R}^{d \times k}$:
{\boldmath
\begin{equation}
W = m \odot \frac{V}{\|V\|_c} = m \odot \frac{W_0 + \Delta W}{\|W_0 + \Delta W\|_c}
\end{equation}
}
where $\|V\|_c = \left[ \|V_{1,:}\|_2, \dots, \|V_{d,:}\|_2 \right]^T \in \mathbb{R}^d$ computes row-wise Euclidean column/row norms, and $\odot$ denotes row-wise scalar multiplication.
Here, $W_0$ is the frozen pretrained weight matrix, $m$ is initialized as $m = \|W_0\|_c$, and $\Delta W$ is parametrized using low-rank matrices:
{\boldmath
\begin{equation}
\Delta W = s \cdot (B \cdot A)
\end{equation}
}
where $A \in \mathbb{R}^{r \times k}$ is Gaussian initialized, $B \in \mathbb{R}^{d \times r}$ is zero-initialized, and $s$ is the scaling factor. In standard LoRA, $s = \frac{\alpha}{r}$; in \textbf{Rank-Stabilized LoRA (rsLoRA)}, $s = \frac{\alpha}{\sqrt{r}}$, which prevents collapse of gradient dynamics as rank $r$ increases.

\subsection{Expanded Non-Technical Definition (Intuition for Anyone)}
\begin{enumerate}
    \item \textbf{Decoupling Volume and Pitch}: Think of tuning a musical instrument. Full fine-tuning adjusts how loud a note is (magnitude) independently from which note is played (direction). Plain LoRA forcibly ties loudness and note pitch together. DoRA separates them completely: you tune volume with a volume knob ($m$) and pitch with a directional compass ($BA$).
    \item \textbf{Closing the Gap with Full Tuning}: Because DoRA mirrors the exact mathematical freedom of full fine-tuning without updating the hundreds of millions of weights in $W_0$, it consistently outperforms standard LoRA and matches full fine-tuning accuracy.
    \item \textbf{Rank-Independent Stability (rsLoRA)}: In standard LoRA, doubling the rank from $r=4$ to $r=64$ shrinks the learning signal by a factor of 16 ($1/r$). rsLoRA divides by $\sqrt{r}$, keeping learning dynamics smooth and scale-invariant across any rank.
    \item \textbf{Zero Inference Penalty}: Once training finishes, the normalized directional matrix $V$ and magnitude vector $m$ are merged into a single deployed dense matrix $W = m \odot \frac{V}{\|V\|_c}$, leaving zero runtime latency.
\end{enumerate}

\subsection{Child-Friendly Mental Model (Physical Metaphor)}
Imagine pushing a sailboat:
\begin{itemize}
    \item \textbf{Standard LoRA}: Tying the rudder angle directly to how hard the engine pushes. If you steer right, you are forced to go faster; if you steer left, you are forced to slow down!
    \item \textbf{DoRA}: Separating the motor throttle ($m$, how strong the push is) from the rudder ($BA$, which direction the boat points). You can push as hard as you want in any direction!
    \item \textbf{Result}: The boat steers exactly like a master captain (full fine-tuning) while only needing two tiny adjustment knobs.
\end{itemize}

\section{Core Mathematical Equations and Definitions}

\subsection{DoRA Forward Weight Formulation}
{\boldmath
\begin{equation}
W = m \odot \frac{W_0 + s \cdot B A}{\|W_0 + s \cdot B A\|_c}
\end{equation}
}
\begin{itemize}
    \item \textbf{Formal Definition}: Parametrization of linear layer weights via explicit magnitude vector $m \in \mathbb{R}^d$ and normalized low-rank directional perturbator.
    \item \textbf{What It Means}: Disentangles scale from orientation for each neuron row.
    \item \textbf{Why It Is Important}: Allows models to perform subtle directional shifts with large magnitude scaling or vice versa, eliminating optimization bottlenecks inherent to additive adapters.
\end{itemize}

\subsection{rsLoRA Scaling Constant}
{\boldmath
\begin{equation}
s_{\text{rsLoRA}} = \frac{\alpha}{\sqrt{r}}
\end{equation}
}
\begin{itemize}
    \item \textbf{Formal Definition}: Rank-stabilized scaling coefficient replacing standard LoRA's $\alpha / r$.
    \item \textbf{What It Means}: Compensates for the fact that singular values of random low-rank products grow as $\sqrt{r}$, preserving learning rate invariance.
    \item \textbf{Why It Is Important}: Guarantees stable convergence for high-capacity adapters ($r = 32, 64, 128$) where standard LoRA stagnates.
\end{itemize}

\subsection{Row-Wise Euclidean Normalization}
{\boldmath
\begin{equation}
\|V_i\|_2 = \sqrt{\sum_{j=1}^k V_{ij}^2 + \epsilon}
\end{equation}
}
\begin{itemize}
    \item \textbf{Formal Definition}: The $L_2$ norm of the $i$-th row of directional matrix $V = W_0 + s BA$.
    \item \textbf{What It Means}: Computes the baseline directional length of neuron weights.
    \item \textbf{Why It Is Important}: Ensures the directional matrix has unit row-norm, isolating all scaling into vector $m$.
\end{itemize}

\section{Rigorous Mathematical Analysis Checklist}

\begin{itemize}
    \item \textbf{Notation}: $W_0 \in \mathbb{R}^{d \times k}$ (frozen base), $m \in \mathbb{R}^d$ (magnitude vector), $A \in \mathbb{R}^{r \times k}, B \in \mathbb{R}^{d \times r}$ (directional low-rank adapter), $\alpha$ (scale), $r$ (rank).
    \item \textbf{Epistemic Status}: Formulated by Liu et al. (ICLR 2024); rsLoRA by Kalajdzievski (2023). Established state of the art in parameter-efficient adaptation.
    \item \textbf{Dependency Graph}: $x \to \text{Evaluate } V = W_0 + sBA \to \text{Row Norm } \|V\|_c \to \text{Normalize } V / \|V\|_c \to \text{Scale by } m \to W x$.
    \item \textbf{Dimensions}: Input $x \in \mathbb{R}^k$, output $y \in \mathbb{R}^d$, rank $r \ll \min(d, k)$.
    \item \textbf{Regimes}: High-performance PEFT, instruction fine-tuning, complex reasoning tasks where standard LoRA underfits.
    \item \textbf{Invariants}: When $B = \mathbf{0}$ and $m = \|W_0\|_c$, $W \equiv W_0$ exactly.
    \item \textbf{Isomorphisms}: Isomorphic to weight normalization (Salimans \& Kingma, 2016) with low-rank directional perturbation.
\end{itemize}

\section{Theoretical Theorems and Formal Proofs}

\begin{theorem}[Initial Equivalence to Pretrained Base Weights]
Let $W_0 \in \mathbb{R}^{d \times k}$ with rows having non-zero norms $\|W_{0, i,:}\|_2 > 0$. If $B = \mathbf{0}_{d \times r}$ and $m_i = \|W_{0, i,:}\|_2$ for all $i \in \{1, \dots, d\}$, then the DoRA weight matrix $W$ is identical to $W_0$:
{\boldmath
\begin{equation}
W_{\text{DoRA}} \equiv W_0
\end{equation}
}
\end{theorem}

\begin{proof}
Let $V = W_0 + s BA$. Since $B = \mathbf{0}_{d \times r}$, for any $A \in \mathbb{R}^{r \times k}$ and $s > 0$:
{\boldmath
\begin{equation}
BA = \mathbf{0}_{d \times k} \implies V = W_0 + \mathbf{0}_{d \times k} = W_0
\end{equation}
}
The row-wise norm of $V$ is therefore:
{\boldmath
\begin{equation}
\|V_{i,:}\|_2 = \|W_{0, i,:}\|_2
\end{equation}
}
Substituting into the DoRA forward formula for row $i$:
{\boldmath
\begin{equation}
W_{i,:} = m_i \frac{V_{i,:}}{\|V_{i,:}\|_2} = \|W_{0, i,:}\|_2 \frac{W_{0, i,:}}{\|W_{0, i,:}\|_2} = W_{0, i,:}
\end{equation}
}
Since this holds for every row $i \in \{1, \dots, d\}$, we have $W \equiv W_0$ identically at step zero. Thus, DoRA preserves the exact pretrained output upon initialization.
\end{proof}

\begin{theorem}[Gradient Scaling Invariance in rsLoRA]
Let $\Delta W = s B A$ where entries of $A$ and $B$ are independent zero-mean random variables with variance $\sigma^2$. If $s = \frac{\alpha}{\sqrt{r}}$, then the expected Frobenius norm $\mathbb{E}[\|\Delta W\|_F^2]$ scales linearly with rank $r$ rather than collapsing as $\mathcal{O}(1/r)$.
\end{theorem}

\begin{proof}
Consider entry $(i, j)$ of product $M = BA$:
{\boldmath
\begin{equation}
M_{ij} = \sum_{l=1}^r B_{il} A_{lj}
\end{equation}
}
Since $B_{il}$ and $A_{lj}$ are independent with mean 0 and variance $\sigma^2$:
{\boldmath
\begin{equation}
\mathbb{E}[M_{ij}] = 0, \quad \operatorname{Var}(M_{ij}) = \sum_{l=1}^r \operatorname{Var}(B_{il}) \operatorname{Var}(A_{lj}) = r \sigma^4
\end{equation}
}
The expected squared Frobenius norm of $M = BA$ is:
{\boldmath
\begin{equation}
\mathbb{E}[\|M\|_F^2] = \sum_{i=1}^d \sum_{j=1}^k \mathbb{E}[M_{ij}^2] = d \cdot k \cdot r \sigma^4
\end{equation}
}
Under standard LoRA scaling $s = \frac{\alpha}{r}$:
{\boldmath
\begin{equation}
\mathbb{E}[\|\Delta W_{\text{LoRA}}\|_F^2] = s^2 \mathbb{E}[\|M\|_F^2] = \frac{\alpha^2}{r^2} (d k r \sigma^4) = \frac{\alpha^2 d k \sigma^4}{r}
\end{equation}
}
As rank $r \to \infty$, the expected norm scales as $\mathcal{O}(1/r) \to 0$, causing gradient vanishing for large ranks.
Under rsLoRA scaling $s = \frac{\alpha}{\sqrt{r}}$:
{\boldmath
\begin{equation}
\mathbb{E}[\|\Delta W_{\text{rsLoRA}}\|_F^2] = \frac{\alpha^2}{r} (d k r \sigma^4) = \alpha^2 d k \sigma^4 = \Theta(1)
\end{equation}
}
The scale of perturbation is independent of rank $r$, establishing rank-stabilized gradient dynamics.
\end{proof}

\section{Foundational Architectural Questions and Epistemic Meta-Questions}

\subsection{Architectural Question 1: Computational Overhead of Normalization}
\textbf{Question:} Does computing row-wise norms in DoRA introduce training slowdown compared to standard LoRA?\\
\textbf{Answer:} During forward and backward propagation, DoRA computes Euclidean row norms and applies chain-rule normalization gradients, introducing approximately $10\%\text{--}15\%$ training runtime overhead. However, memory consumption remains virtually unchanged since only $m$ (of size $d$) and low-rank matrices $A, B$ require optimizer states.

\textbf{Meta-Question:} How is DoRA deployed without incurring latency overhead?\\
\textbf{Answer:} Before exporting to inference engines (e.g. vLLM, TensorRT-LLM), the adapter is fused into a single dense matrix:
{\boldmath
\begin{equation}
W_{\text{fused}} = m \odot \frac{W_0 + s BA}{\|W_0 + s BA\|_c}
\end{equation}
}
Because $W_{\text{fused}}$ has identical dimension $(d, k)$ to $W_0$, inference speed and memory footprint are strictly identical to the unadapted model.

\subsection{Architectural Question 2: Synergy Between DoRA and rsLoRA}
\textbf{Question:} Can rsLoRA scaling be applied directly inside DoRA?\\
\textbf{Answer:} Yes. In DoRA, the directional update $\Delta W$ uses $s \cdot BA$. Setting $s = \alpha / \sqrt{r}$ stabilizes high-rank directional tuning. When scaling rank to $r = 64$ or $r = 128$, rsLoRA prevents directional gradient attenuation, while DoRA magnitude vectors adapt the scale independently.

\textbf{Meta-Question:} Should magnitude vector $m$ be updated with the same optimizer and learning rate as directional matrices $A, B$?\\
\textbf{Answer:} Typically, yes; however, standard practice uses AdamW with weight decay disabled on magnitude parameters $m$ (similar to biases and LayerNorm gains) to avoid shrinkage of neuron magnitudes over time.

\section{Algorithmic Implementation Specification}

\begin{algorithm}[H]
\caption{DoRA Forward Evaluation and Directional Decomposition}
\begin{algorithmic}[1]
\Require Input vector $\mathbf{x} \in \mathbb{R}^k$, base matrix $\mathbf{W}_0 \in \mathbb{R}^{d \times k}$, low-rank matrices $\mathbf{A} \in \mathbb{R}^{r \times k}, \mathbf{B} \in \mathbb{R}^{d \times r}$, magnitude vector $\mathbf{m} \in \mathbb{R}^d$, scaling parameters $\alpha, r$, mode flag $\text{rsLoRA}$
\Ensure Adapted output vector $\mathbf{y} \in \mathbb{R}^d$, directional row norms $\mathbf{n} \in \mathbb{R}^d$, scaling $s$
\If{$\text{rsLoRA}$ is True}
    \State $s \gets \alpha / \sqrt{r}$
\Else
    \State $s \gets \alpha / r$
\EndIf
\State Allocate $\mathbf{y} \in \mathbb{R}^d, \quad \mathbf{n} \in \mathbb{R}^d$
\For{$i = 0$ \textbf{to} $d-1$}
    \State Allocate row $\mathbf{v} \in \mathbb{R}^k, \quad \text{sum\_sq} \gets 0.0$
    \For{$j = 0$ \textbf{to} $k-1$}
        \State $\delta_{ij} \gets \sum_{l=0}^{r-1} \mathbf{B}[i, l] \cdot \mathbf{A}[l, j]$
        \State $\mathbf{v}[j] \gets \mathbf{W}_0[i, j] + s \cdot \delta_{ij}$
        \State $\text{sum\_sq} \gets \text{sum\_sq} + \mathbf{v}[j]^2$
    \EndFor
    \State $\mathbf{n}[i] \gets \sqrt{\max(\text{sum\_sq}, 10^{-12})}$
    \State $\text{dot} \gets 0.0$
    \For{$j = 0$ \textbf{to} $k-1$}
        \State $w_{ij} \gets \mathbf{m}[i] \cdot \left(\mathbf{v}[j] / \mathbf{n}[i]\right)$
        \State $\text{dot} \gets \text{dot} + w_{ij} \cdot \mathbf{x}[j]$
    \EndFor
    \State $\mathbf{y}[i] \gets \text{dot}$
\EndFor
\State \Return $(\mathbf{y}, \mathbf{n}, s)$
\end{algorithmic}
\end{algorithm}

\section{Big-$\Theta$ Complexity Analysis}

\begin{itemize}
    \item \textbf{Directional Reconstruction Time}: $\Theta(d \cdot k \cdot r)$ operations to evaluate $s \cdot B A$ across all elements.
    \item \textbf{Row Norm Evaluation Time}: $\Theta(d \cdot k)$ operations for row-wise sum of squares.
    \item \textbf{Normalized Matrix-Vector Multiplication Time}: $\Theta(d \cdot k)$ multiply-accumulate operations.
    \item \textbf{Total Time Complexity}: $\Theta(d \cdot k \cdot r)$ during training evaluation.
    \item \textbf{Auxiliary Memory Space Complexity}: $\Theta(d)$ space for magnitude vector and directional row norms.
\end{itemize}

\section{Single Canonical Repository Reference}

The complete deterministic scalar implementation, verification suite, and unit test benchmarks are published at:
\begin{center}
\url{https://github.com/Chief-Strategist-J/llm-obs-infra}
\end{center}

\end{document}
